\documentclass[conference]{IEEEtran}
\IEEEoverridecommandlockouts

% --- Encoding & Sprache ---
\usepackage[utf8]{inputenc}
\usepackage[T1]{fontenc}
\usepackage[ngerman]{babel}

% --- Pakete laut IEEE-Template ---
\usepackage{cite}
\usepackage{amsmath,amssymb,amsfonts}
\usepackage{algorithmic}
\usepackage{graphicx}
\usepackage{textcomp}
\usepackage{xcolor}
\usepackage{array}
\usepackage{booktabs}

% --- Inline-Code Befehl ---
\newcommand{\code}[1]{\texttt{\small #1}}

\def\BibTeX{{\rm B\kern-.05em{\sc i\kern-.025em b}\kern-.08em
    T\kern-.1667em\lower.7ex\hbox{E}\kern-.125emX}}

\begin{document}

\title{Container und virtuelle Maschinen --\\
Isolationsmechanismen, Bedrohungsmodelle\\
und Sicherheitsanalyse\\
{\footnotesize IT-Sicherheit Übung -- Thema 12, Hochschule Karlsruhe}}

\author{%
\IEEEauthorblockN{1\textsuperscript{st} Simon Klaes}
\IEEEauthorblockA{\textit{Hochschule Karlsruhe} \\
Karlsruhe, Deutschland \\
klsi1012@h-ka.de}
\and
\IEEEauthorblockN{2\textsuperscript{nd} Radek Pekul}
\IEEEauthorblockA{\textit{Hochschule Karlsruhe} \\
Karlsruhe, Deutschland \\
pera1011@h-ka.de}
\and
\IEEEauthorblockN{3\textsuperscript{rd} Joel Maxeiner}
\IEEEauthorblockA{\textit{Hochschule Karlsruhe} \\
Karlsruhe, Deutschland \\
majo1064@h-ka.de}
\and
\IEEEauthorblockN{4\textsuperscript{th} Pius Tebbe}
\IEEEauthorblockA{\textit{Hochschule Karlsruhe} \\
Karlsruhe, Deutschland \\
tepi1011@h-ka.de}
}

\maketitle

\begin{abstract}
Container und virtuelle Maschinen (VMs) sind die zentralen Bausteine moderner
IT-Infrastrukturen, unterscheiden sich aber grundlegend in ihrer Isolations\-stärke.
Diese Arbeit vergleicht beide Konzepte hinsichtlich ihrer Sicherheits\-eigenschaften:
Sie erläutert die fundamentalen Isolationsmechanismen (Namespaces, cgroups,
Hypervisor-Virtualisierung), demonstriert in einer Live-Demo den Normalbetrieb
beider Technologien sowie typische Fehlkonfigurationen, und diskutiert reale
Escape-Szenarien und Best Practices. Ziel ist es, fundiert zu zeigen, wann welche
Technologie sicherheits\-technisch angemessen ist.
\end{abstract}

\begin{IEEEkeywords}
Container, Docker, virtuelle Maschine, Hypervisor, Namespaces, cgroups,
Isolation, Container-Escape, IT-Sicherheit
\end{IEEEkeywords}

% =====================================================================
% Person A -- ca. 1 Seite (ca. 1 Spalte in IEEE Conference Format)
% =====================================================================
\section{Grundlagen der Isolation}

\subsection{Motivation und Bedrohungsmodell}
Sobald mehrere Workloads auf derselben Hardware laufen, stellt sich die Frage,
\textit{welche Komponente bei einem Ausbruch versagen darf}. Beim
\textbf{Container-Bedrohungsmodell} wird angenommen, dass der Angreifer bereits
Code im Container ausführt -- etwa über eine kompromittierte Webanwendung
oder ein bösartiges Image -- und versucht, den Host oder Nachbarcontainer zu
übernehmen. Beim \textbf{VM-Bedrohungsmodell} wird sogar angenommen, dass der
Gast vollständig kompromittiert ist (Root im Gast-Kernel); dennoch sollen Host
und andere Gäste sicher bleiben. Container schützen damit \textit{vor versehentlichen
Konflikten zwischen vertrauenswürdigen Workloads}, VMs dagegen
\textit{vor bösartigen Mandanten} und sind deshalb Standard für Multi-Tenant-Clouds
wie AWS EC2 oder Azure VMs~\cite{b_lehrgang}.

\subsection{Virtualisierungskonzepte: Hypervisor-Typen}
Eine VM simuliert einen kompletten Computer mit eigenem Kernel, eigener
Hardware-Sicht und eigenem Bootloader. Der Hypervisor ist die Schicht, die
diese Trennung erzwingt:
\begin{itemize}
  \item \textbf{Typ 1 (Bare-Metal):} läuft direkt auf der Hardware, ohne Host-OS
        dazwischen. Beispiele: KVM, Xen, VMware ESXi, Microsoft Hyper-V.
        Geringe Trusted Computing Base (TCB) und produktiv bevorzugt.
  \item \textbf{Typ 2 (Hosted):} läuft als Anwendung auf einem Betriebssystem.
        Beispiele: VirtualBox, VMware Workstation, Parallels Desktop, UTM.
\end{itemize}
Hardware-Erweiterungen (Intel VT-x, AMD-V, ARM Virtualization Extensions)
führen einen privilegierteren CPU-Modus ein (\glqq Ring~$-1$\grqq{} bzw.
EL2 auf ARM). Privilegierte Operationen des Gasts lösen einen
\textit{VM Exit} aus, den der Hypervisor behandelt. EPT/NPT übersetzt
Gast-physikalische zu Host-physikalischen Adressen in Hardware. Die
Schnittstelle zwischen Gast und Host besteht damit nur aus wenigen Hypercalls
und para\-virtualisierten Geräten (virtio) -- eine deutlich kleinere
Angriffsfläche als das Linux-Syscall-Interface~\cite{b_bedrohung}.

\subsection{Container-Grundlagen: Namespaces, Cgroups, Kernel-Sharing}
Container sind technisch \textbf{keine eigene Technologie}, sondern eine
Kombination bestehender Linux-Kernel-Features. Sie sind im Wesentlichen
isolierte Prozesse, die denselben Kernel mit dem Host teilen.

\textbf{Namespaces} schaffen pro Container eine eigene Sicht auf Kernel-Ressourcen:
\code{PID} (eigene Prozessbäume), \code{NET} (eigener Netzwerkstack), \code{MNT}
(eigene Mount-Points), \code{IPC}, \code{UTS} (Hostname), \code{USER} (UID-Mapping),
\code{CGROUP} und \code{TIME}. Namespaces beschränken \textit{Sichtbarkeit}, nicht
welche Syscalls ausgeführt werden dürfen.

\textbf{Control Groups (cgroups)} sind orthogonal dazu und beschränken
\textit{Ressourcenverbrauch} (CPU, Speicher, Block-I/O, Anzahl Prozesse). Ohne
cgroups könnte ein einzelner Container per Fork-Bombe das Host-System lahm\-legen.

Ergänzend reduzieren weitere Mechanismen die Angriffsfläche:
\textbf{Capabilities} zerlegen die Root-Rechte in ca. 40 Einzelrechte (Container
erhalten standardmäßig nur ca. 14), \textbf{Seccomp-BPF} filtert Syscalls (Docker
blockiert per Default ca. 60 von ca. 330), und \textbf{LSMs} wie AppArmor oder
SELinux setzen Mandatory Access Control durch.

\subsection{Vertrauensgrenzen und Vergleich}
Der entscheidende Unterschied ist die Frage, was bei einem Ausbruch versagen
muss. Bei Containern ist die Vertrauensgrenze der \textbf{geteilte Linux-Kernel}
-- jede ausnutzbare Kernel-Schwachstelle ist potenziell ein Container-Escape.
Die TCB umfasst ca. 30 Millionen Zeilen Code (gesamter Kernel inkl. Treiber,
Syscall-Interface, Dateisystemen). Bei VMs ist die Vertrauensgrenze der
\textbf{Hypervisor}; die TCB im Hot-Path liegt z.\,B. bei KVM bei rund
100\,000 Zeilen Code -- zwei bis drei Größenordnungen weniger
Angriffsfläche~\cite{b_lehrgang}. Tabelle~\ref{tab:isolation} fasst die
wesentlichen Dimensionen zusammen.

\begin{table}[htbp]
\caption{Vergleich der Isolationsgarantien}
\label{tab:isolation}
\begin{center}
\footnotesize
\begin{tabular}{@{}p{2.0cm}p{2.4cm}p{2.6cm}@{}}
\toprule
\textbf{Dimension} & \textbf{Container} & \textbf{VM} \\
\midrule
Isolationsgrenze & Kernel (geteilt)        & Hardware (Hypervisor) \\
Boot-Zeit        & Millisekunden           & Sekunden bis Minuten \\
Overhead         & MB                      & Hunderte MB bis GB \\
Dichte/Host      & Hunderte                & Dutzende \\
Angriffsfläche   & $\sim$330 Syscalls      & Hypercalls + Geräte \\
TCB              & $\sim$30 Mio. LOC       & $\sim$100\,k LOC \\
Multi-Tenancy    & Unzureichend            & Standard \\
\bottomrule
\end{tabular}
\end{center}
\end{table}

% =====================================================================
% Person B -- ca. 1 Seite -- Platzhalter
% =====================================================================
\section{Implementierung und Normalbetrieb}
\textit{[Abschnitt von Person B: Setup des Demo-Dienstes (Python HTTP-Server)
sowohl als Docker-Container als auch in einer VirtualBox/UTM-VM,
Ressourcenisolierung mit \code{docker stats} und \code{htop}, Netzwerkvergleich
Container-Bridge vs.\ NAT/Bridged, Analyse von \code{ps}, \code{top},
\code{/proc}, Fazit zur vorgesehenen Isolation.]}

% =====================================================================
% Person C -- ca. 1 Seite -- Platzhalter
% =====================================================================
\section{Fehlkonfigurationen und Schwachstellen}
\textit{[Abschnitt von Person C: Demonstration und Analyse typischer
Fehlkonfigurationen. Container: \code{--privileged}, gemountete Host-Volumes
(\code{-v /etc:/host-etc}), deaktiviertes Seccomp-Profil
(\code{--security-opt seccomp=unconfined}), gemounteter Docker-Socket. VM:
unsichere Bridged-Konfiguration, Shared Folders, aktive Komfort-Geräte.
Für jede Fehlkonfiguration: welche Isolationsgrenze bricht?]}

% =====================================================================
% Person D -- ca. 1 Seite -- Platzhalter
% =====================================================================
\section{Escape-Szenarien und Best Practices}
\textit{[Abschnitt von Person D: Container-Escapes mit CVE-Beispielen
(CVE-2019-5736 runc, CVE-2022-0492 cgroups v1, Dirty Pipe CVE-2022-0847),
VM-Escapes (VENOM CVE-2015-3456, VirtualBox-CVEs), Side-Channel-Angriffe
(Spectre/Meltdown). Best Practices: Rootless Container, minimale Images,
seccomp, read-only Filesystem, Netzwerksegmentierung, Hypervisor-Härtung.
Gesamtfazit: wann Container, wann VM?]}

% =====================================================================
\section*{Hinweis zur KI-Nutzung}
Bei der Recherche und der Erstellung dieses Dokuments wurden KI-gestützte
Werkzeuge (Anthropic Claude) zur Strukturierung von Inhalten, zum Erklären
von Fachbegriffen und zur sprachlichen Glättung eingesetzt. Sämtliche
inhaltlichen Aussagen wurden von den Autorinnen und Autoren geprüft und
gegen die Lehrgangs- und Referenzquellen abgeglichen.

% =====================================================================
\begin{thebibliography}{00}
\bibitem{b_lehrgang} Systemsicherheits-Lehrgang, \glqq Container vs.\ Virtuelle Maschinen -- Ein tiefgehender Lehrgang\grqq, Hochschule Karlsruhe, 2026.
\bibitem{b_bedrohung} \glqq Container und Virtuelle Maschinen -- Isolationsmechanismen, Bedrohungsmodelle und Sicherheitsanalyse\grqq, internes Skript, Hochschule Karlsruhe, 2026.
\bibitem{b_cve5736} MITRE, CVE-2019-5736 (\code{runc}-Container-Breakout), 2019.
\bibitem{b_dirtypipe} M.\ Kellermann, \glqq The Dirty Pipe Vulnerability (CVE-2022-0847)\grqq, 2022.
\bibitem{b_venom} CrowdStrike, \glqq VENOM Vulnerability (CVE-2015-3456)\grqq, 2015.
\end{thebibliography}

\end{document}
